\documentclass[10pt]{article}

% --- Encoding & fonts -------------------------------------------------------
\usepackage[utf8]{inputenc}   % accented author names (Tönnis, Reikerås, ...)
\usepackage[T1]{fontenc}
\usepackage{textcomp}         % \textdegree
\usepackage{amsmath}          % \text, \tfrac
\usepackage{lmodern}

% --- Page geometry: landscape to fit the five wide columns ------------------
\usepackage[a4paper,landscape,margin=1.2cm]{geometry}

% --- Tables -----------------------------------------------------------------
\usepackage{array}
\usepackage{booktabs}
\usepackage{longtable}
\usepackage{enumitem}
\usepackage{graphicx}         % illustration figures (vector PDFs)
\graphicspath{{illustrations/pdf/}}
\usepackage{xcolor}
\definecolor{refblue}{HTML}{2563EB}
\definecolor{refred}{HTML}{DC2626}
\definecolor{refgreen}{HTML}{0E9F6E}
\definecolor{refpurple}{HTML}{7C3AED}
\newcommand{\incl}[1]{\includegraphics[width=8.4cm,height=6.4cm,keepaspectratio]{#1}}

% Left-aligned, ragged-right fixed-width paragraph column.
\newcolumntype{L}[1]{>{\raggedright\arraybackslash}p{#1}}

\setlength{\tabcolsep}{4pt}
\renewcommand{\arraystretch}{1.15}

\title{Morphometric Parameter Reference Table}
\date{}

\begin{document}
\maketitle

A companion to \texttt{reader\_measurement\_guide.md}. For each parameter measured by
the pipeline this table gives the textbook definition, the clinical rationale, a
step-by-step summary of how the algorithm implements it (condensed from the guide),
and a normal-adult reference range with a literature citation.

\subsection*{Important caveats}
\begin{itemize}[leftmargin=1.4em,itemsep=2pt,topsep=2pt]
  \item Reference ranges are \textbf{method- and modality-dependent}. Where the
        algorithm's landmark choice materially shifts the expected value (most
        strongly for femoral and tibial torsion), this is flagged in the range
        column. Always interpret the algorithm's output against a range obtained
        with a comparable technique.
  \item Distances are in millimetres in patient (world) coordinates; angles are
        unsigned degrees unless a sign convention is stated.
\end{itemize}

Citations are keyed [n] to the References section.

\begin{footnotesize}
\begin{longtable}{L{2.5cm} L{4.0cm} L{4.0cm} L{9.2cm} L{5.4cm}}
\toprule
\textbf{Parameter} & \textbf{Textbook definition} & \textbf{Clinical rationale} &
\textbf{Implementation} & \textbf{Reference range} \\
\midrule
\endfirsthead

\multicolumn{5}{l}{\emph{\ldots continued}} \\
\toprule
\textbf{Parameter} & \textbf{Textbook definition} & \textbf{Clinical rationale} &
\textbf{Implementation} & \textbf{Reference range} \\
\midrule
\endhead

\midrule
\multicolumn{5}{r}{\emph{continued on next page\ldots}} \\
\endfoot

\bottomrule
\endlastfoot

% --- 1. Femoral torsion ----------------------------------------------------
\textbf{Femoral torsion (anteversion)} &
Transverse-plane angle between the femoral neck axis (proximal reference) and the
posterior condylar line of the distal femur (distal reference). &
Quantifies rotational deformity of the femur; excessive anteversion or retroversion
contributes to patellofemoral instability, gait abnormality, in-/out-toeing, and
femoroacetabular impingement, and guides derotation osteotomy planning. &
(1) Fit a sphere to the femoral head $\rightarrow$ centre and radius \emph{r}.
(2) \textbf{Proximal line, two methods.} \emph{Lee:} on the first inferior slice
where the femur is a single outline and a circle of radius $2\cdot r$ still hits
bone, keep outline points in a ring ${\approx}0.9\text{--}1.1\cdot r$ near the neck
centre of mass and fit a line through them and the head centre. \emph{Murphy:} find
the slice where the lesser trochanter is most prominent; the neck base is the shaft
cross-section centroid there; the line runs head centre $\rightarrow$ neck base.
(3) \textbf{Distal line:} on the slice with the largest convex condylar bone area,
the tangent to the two posterior condyles. (4) Measure each line's angle to the
medial--lateral axis in the transverse plane, fold to 0--90\textdegree{}, and add
(opposite AP tilt) or subtract (same tilt). Both Lee and Murphy values are reported. &
Strongly method-dependent. \emph{Lee} (proximal neck landmark) $\approx$
\textbf{8--15\textdegree{}} (CT 11\textdegree{}$\pm$11\textdegree{}); \emph{Murphy}
(distal neck landmark) $\approx$ \textbf{18--28\textdegree{}}
(CT 28\textdegree{}$\pm$13\textdegree{}) --- the two differ by ${\sim}17$\textdegree{}
on the same hip. Classic mixed-method normal $\approx$ 10--15\textdegree{}
(broad 8--20\textdegree{}); retroversion $<$5\textdegree{}, excessive $>$20\textdegree{}.
CT and MRI agree closely. [1][2] \\
\midrule

% --- 2. Tibial torsion -----------------------------------------------------
\textbf{Tibial torsion} &
Transverse-plane angle between the posterior condylar line of the proximal tibia and
the distal trans-malleolar / tibia--fibula line. &
Quantifies rotational deformity of the tibia; excessive external torsion causes
out-toeing and patellofemoral malalignment; relevant to derotation osteotomy and to
interpreting overall lower-limb torsional profile. &
(1) \textbf{Proximal line:} as for the femoral condylar line (\S1.2) but on the tibia
--- on the proximal tibial slice with the largest convex bone area, the tangent to the
two posterior plateau aspects. (2) \textbf{Distal line:} on the slice with the largest
tibial (equivalent-ellipse) diameter at the ankle, connect the tibia centroid to the
fibula centroid. (3) Measure each line's angle to the medial--lateral axis, fold to
0--90\textdegree{}, add/subtract by AP tilt direction. &
External (positive) and governed by the \emph{distal} landmark: trans-malleolar axis
$\approx$ \textbf{40\textdegree{}$\pm$9\textdegree{}}; bimalleolar axis $\approx$
\textbf{20--30\textdegree{}}. Whole-limb CT cohorts $\approx$ 33--35\textdegree{}. This
pipeline uses tibia--fibula centroids (closer to the trans-malleolar family). MRI is a
valid CT substitute. [3][4] \\
\midrule

% --- 3. Knee rotation angle ------------------------------------------------
\textbf{Knee rotation angle} &
Transverse-plane angle between the posterior condylar line of the distal femur and the
posterior condylar line of the proximal tibia (rotational femoro-tibial mismatch). &
Captures rotational alignment of the knee joint itself; relevant to patellofemoral
tracking and to component rotation in arthroplasty planning. &
(1) \textbf{Femoral line} as in femoral torsion \S1.2. (2) \textbf{Tibial line} as in
tibial torsion \S2.1, both taken at the knee. (3) Transverse-plane angle between them,
folded to 0--90\textdegree{} and combined by AP tilt direction; an exact 180\textdegree{}
result is reset to 0\textdegree{}. &
No standardized normative range for this exact bony pairing exists in the literature;
the posterior condylar lines are near-parallel in extension, so the isolated angle is
small (single-digit degrees). Nearest published surrogates (different reference lines):
EOS femoro-tibial rotation 2.3\textdegree{}$\pm$7.2\textdegree{}; CT tibiofemoral
rotation ${\sim}4.1$\textdegree{}$\pm$5.3\textdegree{}. State your reference lines
explicitly. [5][6] \\
\midrule

% --- 4. CCD angle ----------------------------------------------------------
\textbf{CCD angle (caput--collum--diaphyseal / neck--shaft)} &
Angle between the femoral shaft axis and the femoral neck axis. &
Used to detect deformities such as coxa valga and coxa vara, and to plan hip
replacement or osteosynthesis after proximal femur fractures. &
(1) Fit sphere to femoral head to determine its centre. (2) Define femoral neck centre
as the centroid of the femoral voxels in a hollow shell (between \emph{r} and
$1.2\cdot r$) that lie both distal and lateral to the head centre; neck axis = head
centre $\rightarrow$ neck centre. (3) Determine the shaft axis from the distal-slice
centroid of the proximal femur to the proximal-slice centroid of the knee femur (or,
without a knee image, up to the greater-trochanter region). (4) Compute the angle
between shaft and neck axes, reported as the obtuse value; also a coronal-plane
projection (AP component zeroed). &
Adult normal $\approx$ \textbf{126--132\textdegree{}} (rotation-corrected CT:
M 129.6\textdegree{}$\pm$5.9\textdegree{}, F 131.9\textdegree{}$\pm$6.8\textdegree{});
teaching range 120--135\textdegree{}. Coxa vara $<$120\textdegree{}, coxa valga
$>$135\textdegree{}. Reference plane and rotation correction shift the value by a few
degrees. [7] \\
\midrule

% --- 5. Bone length --------------------------------------------------------
\textbf{Bone length (femur / tibia / whole leg)} &
Straight-line (Euclidean) distance in mm between the proximal-most and distal-most
bone centroids. &
Quantifies limb and segment length for leg-length-discrepancy assessment and
surgical/orthotic planning. &
(1) \textbf{Proximal point:} centre of mass on the most proximal (superior) bone slice.
(2) \textbf{Distal point:} centroid of the most distal slice (femur / whole leg), or
--- for the tibia --- the centroid of the distal articulating surface (plafond slice,
the most distal slice before the cross-sectional area changes abruptly), not the
malleolus tip. (3) 3D distance between the two centroids in patient coordinates (mm),
accounting for slice spacing. &
Highly sex- and population-dependent. Radiographic (standing, full-length): femur
$\approx$ \textbf{500 mm} (range ${\sim}393\text{--}584$), tibia $\approx$
\textbf{390 mm} (range ${\sim}308\text{--}465$), femorotibial ratio ${\sim}1.28{:}1$.
Cadaveric ``maximum length'' runs ${\sim}5\text{--}6$ cm shorter than radiographic
length --- definitions are not interchangeable. [8] \\
\midrule

% --- 6. Acetabular version -------------------------------------------------
\textbf{Acetabular version (anteversion)} &
Transverse-plane angle describing how anteverted the acetabular opening is, on a single
transverse slice. &
Abnormal acetabular version (retroversion or excessive anteversion) is associated with
femoroacetabular impingement, instability, and hip osteoarthritis. &
(1) Locate both femoral head centres (sphere fit); the measurement slice is midway
between them. (2) On that slice, per acetabulum, find the posterior rim point p1 (most
lateral bone in the posterior portion) and anterior rim point p2 (most lateral bone in
the anterior portion). (3) Build reference line G joining the two head centres.
(4) Construct the perpendicular to G through p1; the version is the angle between the
rim line (p1 $\rightarrow$ p2) and that perpendicular. &
$\approx$ \textbf{15\textdegree{}$\pm$6\textdegree{}}; normal band roughly
\textbf{10--25\textdegree{}} ($<$10\textdegree{} too little, $>$25\textdegree{} too much,
both linked to OA). CT cohort means cluster 16--21\textdegree{}, a few degrees higher in
women. Strongly plane-/modality-dependent. [9] \\
\midrule

% --- 7. Center-edge angle --------------------------------------------------
\textbf{Center-edge (CE) angle of Wiberg (LCEA)} &
Coronal-plane angle between a vertical reference through the femoral head centre and the
line from the head centre to the lateral edge of the acetabular roof. &
Primary measure of lateral acetabular coverage of the femoral head; low values indicate
hip dysplasia (undercoverage), high values pincer-type overcoverage. &
(1) Locate both femoral head centres (sphere fit) and reference line G joining them; the
vertical reference is the line through the head centre perpendicular to G, directed
superiorly. (2) Lateral acetabular edge = the most lateral (then most superior) bony
point of the acetabular roof above the head (search starts ${\sim}1.1\text{--}1.5\cdot r$
above the head centre). (3) CE angle = angle between the vertical reference and the
head-centre $\rightarrow$ lateral-edge line, in the coronal plane. &
Mean \textbf{33.6\textdegree{}} (2.5th--97.5th pct 18.1--48.0\textdegree{}). Thresholds:
$<$20\textdegree{} dysplasia, 20--25\textdegree{} borderline, \textbf{25--40\textdegree{}
normal}, $>$40\textdegree{} overcoverage (pincer FAI). Underestimated on plain films vs
3D CT. [10] \\
\midrule

% --- 8. Femoral offset -----------------------------------------------------
\textbf{Femoral offset} &
Perpendicular distance (mm) from the femoral head centre to the femoral shaft axis. &
Determines the hip abductor moment arm; restoring native offset in total hip
arthroplasty is important for abductor strength, stability, and wear. &
(1) Femoral head centre (sphere fit). (2) Femoral shaft axis (same construction as CCD
\S4.1 step 3). (3) Drop a perpendicular from the head centre onto the shaft axis;
offset = its length in mm in patient coordinates. (A projected variant zeroes the AP
component for a pure coronal measurement; default is full 3D.) &
Mean $\approx$ \textbf{43 mm} (43.0$\pm$6.8 mm; range ${\sim}24\text{--}61$ mm). Larger
in men, scales with femur size; plain radiographs underestimate vs CT. [11] \\
\midrule

% --- 9. HKA ----------------------------------------------------------------
\textbf{Hip--knee--ankle (HKA) angle} &
Coronal-plane mechanical-axis angle at the knee between the femoral-head$\rightarrow$knee
segment and the knee$\rightarrow$ankle segment. &
The principal measure of coronal limb alignment (varus/valgus); central to osteotomy and
arthroplasty planning and to osteoarthritis risk assessment. &
(1) \textbf{Hip point:} femoral head centre (sphere fit). (2) \textbf{Knee point:} tibia
centre of mass in the knee region. (3) \textbf{Ankle point:} centroid of the distal
tibial articulating surface (plafond slice). (4) Project hip$\rightarrow$knee and
knee$\rightarrow$ankle vectors to the coronal plane (AP component zeroed) and report the
angle between them; a straight axis $\approx$ 0\textdegree{}, varus/valgus increases it. &
Neutral $\approx$ \textbf{0\textdegree{}} deviation; mechanically-normal limbs run to
${\sim}\textbf{1.5\textdegree{} varus}$ (constitutional aHKA
1.14\textdegree{}$\pm$1.85\textdegree{} varus; normal band ${\sim}1$\textdegree{}
valgus--3.2\textdegree{} varus). Equivalent to conventional HKA $\approx$
178.5\textdegree{}$\pm$2\textdegree{} (180\textdegree{} = straight; $<$180\textdegree{}
varus, $>$180\textdegree{} valgus). General-population SD is wider (${\sim}\pm5$\textdegree{}).
[12][13] \\
\midrule

% --- 10. JLCA --------------------------------------------------------------
\textbf{Joint-line convergence angle (JLCA)} &
Coronal-plane angle between the distal femoral joint line and the proximal tibial joint
line. &
Reflects intra-articular (cartilage/soft-tissue) contribution to coronal deformity;
distinguishes bony from intra-articular malalignment in osteotomy planning. &
(1) \textbf{Femoral joint line:} from transverse scanning, the two distal-most condylar
centres (smaller-condyle centroid on the most distal two-component slice; the other from
the most distal single-component slice), projected to the coronal plane. (2) \textbf{Tibial
joint line:} intercondylar eminence = centroid of the most proximal tibial slice; restrict
to the outer thirds of each plateau (lateral/medial of the eminence by $>\tfrac{2}{3}$ of
the eminence-to-edge distance), take the most proximal point in each region (the plateau
apices), projected to the coronal plane. (3) JLCA = angle between the two lines, folded to
$\leq$90\textdegree{}. &
$\approx$ \textbf{0--2\textdegree{}} (joint lines near-parallel); healthy cohort
1.2\textdegree{}$\pm$2.7\textdegree{}. Lateral opening conventionally positive; magnitude
rises with OA severity and laxity. [13][14] \\
\midrule

% --- 11. MAD ---------------------------------------------------------------
\textbf{Mechanical-axis deviation (MAD)} &
Perpendicular distance (mm) from the knee centre to the mechanical (Mikulicz) axis, in
the coronal plane. &
Quantifies the magnitude of coronal malalignment in millimetres; used alongside HKA to
plan and verify deformity correction. &
(1) Femoral head centre (sphere fit). (2) Ankle centre = centroid of the distal tibial
articulating surface (plafond slice). (3) Mikulicz line = head centre $\rightarrow$ ankle
centre. (4) Knee centre = tibia centre of mass in the knee region. (5) Project all three
to the coronal plane; MAD = perpendicular distance from the knee centre to the Mikulicz
line. \textbf{Sign:} negative for medial deviation, positive for lateral. &
Mikulicz line normally passes ${\sim}\textbf{8 mm$\pm$7 mm medial}$ to the knee centre;
normal commonly stated as \textbf{1--15 mm medial} ($>$15 mm medial = varus/abnormal).
Alternative refs cite 4$\pm$2 mm medial. Physiological normal is a small medial offset,
not zero. [15] \\
\end{longtable}
\end{footnotesize}

\clearpage
\section*{Illustrations}

Schematic diagrams of each parameter --- the reference lines, landmarks, and the
measured angle / distance. They are didactic, \emph{not} anatomically exact, and
not for clinical use. Colour key:
\textcolor{refblue}{\textbf{blue}} = proximal / first reference,
\textcolor{refred}{\textbf{red}} = distal / second reference,
grey dashed = construction,
\textcolor{refgreen}{\textbf{emerald}} = measured angle,
\textcolor{refpurple}{\textbf{purple}} = measured distance. The JLCA and
knee-rotation angles are drawn deliberately exaggerated (their true values are
near-parallel).

% One tabular per row (each figure self-labels via its own title) so the grid
% can break across pages — a single tabular cannot.
\newcommand{\illrow}[1]{\begin{center}\setlength{\tabcolsep}{10pt}\begin{tabular}{ccc}#1\end{tabular}\end{center}\vspace{6pt}}

\illrow{\incl{femoral_torsion} & \incl{tibial_torsion} & \incl{knee_rotation}}
\illrow{\incl{ccd_angle} & \incl{bone_length} & \incl{acetabular_version}}
\illrow{\incl{ce_angle} & \incl{femoral_offset} & \incl{hka_angle}}
\illrow{\incl{jlca} & \incl{mad} &}

\section*{References}
\begin{enumerate}[leftmargin=2em,itemsep=2pt]
  \item Schmaranzer F, Lerch TD, Siebenrock KA, et al. Differences in Femoral Torsion
        Among Various Measurement Methods Increase in Hips With Excessive Femoral
        Torsion. \emph{Clin Orthop Relat Res.} 2019;477(5):1073--1083.
  \item Murphy SB, Simon SR, Kijewski PK, et al. Femoral anteversion. \emph{J Bone Joint
        Surg Am.} 1987;69(8):1169--1176.
  \item Jend HH, et al. CT determination of tibial torsion (trans-malleolar axis method).
        1981. (See also Strecker W, Keppler P, Gebhard F, Kinzl L. Length and torsion of
        the lower limb. \emph{J Bone Joint Surg Br.} 1997;79-B(6):1019--1023.)
  \item Schneider B, Laubenberger J, et al. Measurement of femoral antetorsion and tibial
        torsion by MR imaging. \emph{Br J Radiol.} 1997;70:575--579.
  \item Than P, Szuper K, Somoskeöy S, et al. Geometrical values of the normal and
        arthritic hip and knee detected with the EOS imaging system. \emph{Int Orthop.}
        2012;36(6):1291--1297.
  \item Lu Y, Ren X, Liu B, et al. Tibiofemoral rotation alignment in normal knee joints
        among Chinese adults: a CT analysis. \emph{BMC Musculoskelet Disord.} 2020;21:323.
  \item Boese CK, Jostmeier J, Oppermann J, et al. The neck-shaft angle: CT reference
        values of 800 adult hips. \emph{Skeletal Radiol.} 2016;45(4):455--463. (See also
        the radiographic systematic review, \emph{Skeletal Radiol.} 2016;45(1):19--28,
        pooled 128.8\textdegree{}.)
  \item Aitken SA. Normative Values for Femoral Length, Tibial Length, and the
        Femorotibial Ratio in Adults Using Standing Full-Length Radiography.
        \emph{Osteology.} 2021;1(2):86--95.
  \item Tönnis D, Heinecke A. Acetabular and femoral anteversion: relationship with
        osteoarthritis of the hip. \emph{J Bone Joint Surg Am.} 1999;81(12):1747--1770.
        (Foundational CT method: Reikerås O, Bjerkreim I, Kolbenstvedt A. \emph{Acta Orthop
        Scand.} 1983;54(1):18--23.)
  \item Werner CMC, Ramseier LE, Ruckstuhl T, et al. Normal values of Wiberg's lateral
        center-edge angle and Lequesne's acetabular index --- a coxometric update.
        \emph{Skeletal Radiol.} 2012;41(11):1273--1282. (Original: Wiberg G. \emph{Acta Chir
        Scand.} 1939.)
  \item Noble PC, Alexander JW, Lindahl LJ, et al. The anatomic basis of femoral component
        design. \emph{Clin Orthop Relat Res.} 1988;(235):148--165.
  \item Lan et al. Normative distribution of the arithmetic hip--knee--ankle angle in
        individuals with normal mechanical alignment. \emph{The Knee}, 2024 (aHKA
        1.14\textdegree{}$\pm$1.85\textdegree{} varus; normal band 1\textdegree{}
        valgus--3.2\textdegree{} varus).
  \item Kura R, et al. Normative values of radiographic parameters in coronal-plane
        lower-limb alignment in a general Japanese population (Iwaki cohort). 2025 (HKA
        1.8\textdegree{}$\pm$5.1\textdegree{} varus; JLCA 1.2\textdegree{}$\pm$2.7\textdegree{}).
  \item Managing intra-articular deformity in high tibial osteotomy: a narrative review.
        \emph{J Exp Orthop.} 2020;7:64 (defines JLCA, normal ${\sim}0\text{--}2$\textdegree{},
        lateral-opening-positive).
  \item Paley D. \emph{Principles of Deformity Correction.} Springer, 2002 (MAD
        ${\sim}8$ mm$\pm$7 mm medial; normal 1--15 mm medial; medial = varus, lateral =
        valgus). See also Sabharwal S, et al. Radiological assessment of lower-limb
        alignment. \emph{EFORT Open Rev.} 2021;6(6):487--494.
\end{enumerate}

\end{document}
